% Source: https://artofproblemsolving.com/community/c3h1566650_mathcount_20172018_handbook_warmup_13____q155_158159
% Stable identity: aops-topic-1566650-q159
Write $A,B,M,N$ for the four numbers: $A\in\{2,3,5,7\}$, $N\in\{1,4,9\}$, all four distinct. Since the numbers are distinct, $2B$ equals exactly one of $A$, $M$, $N$; we split into these three cases, which do not overlap.

\textbf{Case $2B=A$.} $A$ is an even prime, so $A=2$, $B=1$ and $N\in\{4,9\}$. $M$ is one of $A+B=3$, $A+N$, $B+N$, at most $10$ and different from the others:
\begin{itemize}
\item $N=4$: $M\in\{3,6,5\}$, three assignments;
\item $N=9$: $M\in\{3,10\}$ ($11$ is too big), two assignments.
\end{itemize}
\textbf{Case $2B=N$.} $N$ is an even square, so $N=4$, $B=2$ and $A\in\{3,5,7\}$. $M$ is one of $A+2$, $A+4$, $6$:
\begin{itemize}
\item $A=3$: $M\in\{5,7,6\}$; $A=5$: $M\in\{7,9,6\}$; $A=7$: $M\in\{9,6\}$ ($11$ is too big),
\end{itemize}
eight assignments.

\textbf{Case $2B=M$.} $M=A+B$ or $M=B+N$ would force $B=A$ or $B=N$, so $M=A+N=2B$. Then $A+N$ is even and at most $10$: $N=1$ with $A\in\{3,5,7\}$, giving $(A,B,M,N)=(3,2,4,1),(5,3,6,1),(7,4,8,1)$; or $N=4$ with $A=2$, giving $(2,3,6,4)$; $N=9$ gives $A+N\ge12$. Four assignments, all with distinct numbers.

In total there are $5+8+4=\boxed{17}$ assignments.
